% Avsnitt 7.7, ur lectures/L07/appendix/b_exercises.md.

\section{Övningar}
\label{c7:sec:exercises}

\begin{aside}
\textbf{Så kontrollerar du ditt arbete.} Varje övning nedan som ber dig skriva en VHDL-modul kommer
med en självkontrollerande testbänk under \file{lectures/L07/exercises/}. Skriv din modul i dess
katalog \file{exercises/<module>/}, med det entitetsnamn och den \textbf{portordning} övningen
anger, och kör den sedan med GHDL \textemdash{} se \secref{c2:sec:testbenches}.

Från och med det här kapitlet återanvänder designerna moduler du skrivit tidigare, och de delas inte
ut en andra gång: varje övning säger vilka av dina egna filer som ska kopieras in, och skriver ut
den \code{ghdl -a}-rad som behöver dem.
\end{aside}

% ------------------------------------------------------------------------------------------
\exerciseset{Timers för hand}

\exercise{En timer med ett annat målvärde}{Krets}
\label{c7:ex:byhand}
Bygg en andra timer för hand i CircuitVerse, enligt \secref{c7:sec:byhand}, men med målvärdet
\code{TICK_COUNT = 12} i stället för avsnittets \code{10}.

\xpart{a} Rikta om jämföraren:
\begin{itemize}
  \item Skriv \code{12} binärt med 4 bitar.
  \item De fyra XNOR-grindarna sitter kvar precis där de sitter. Säg vilka av deras konstanta
    ingångar som ändras, och till vad.
  \item Säg, i en mening, varför en likhetsjämförare byggd av XNOR inte behöver någon grind
    tillagd, borttagen eller omkopplad för att jämföra mot ett annat värde.
\end{itemize}

\xpart{b} Bygg den: återanvänd din 4-bitars räknare och jämförare från passet, och ändra bara
konstanterna. Sätt klockperioden till \code{1000 ms} och bekräfta att \code{timeout} slår till en
gång var trettonde klockflank, under exakt en flank varje gång.

\xpart{c} Ta nu bort nollställningen, så att \code{timeout} inte längre tvingar räknaren tillbaka
till \code{0000}. Förutsäg, innan du kör den, hur ofta \code{timeout} slår till nu. Kör den och
bekräfta. Förklara resultatet utifrån räknarens eget överslag (\secref{c6:sec:counter}).

\xpart{d} Gör sönder jämföraren med flit: koppla loss en XNOR:s utgång från AND-grinden och håll den
AND-ingången hög i stället. Förutsäg vilka räknarvärden som nu höjer \code{timeout}, innan du
simulerar, och bekräfta genom att stega räknaren genom en hel cykel. Det här är hur ”tre av de fyra
bitarna stämmer” ser ut, och det är det vanligaste sättet för en handritad jämförare att bli fel.

\note[Tips]Att räkna \code{0} till och med \code{12} är tretton flanker, inte tolv.
\Secref{c7:sec:concept}:s anmärkning om \code{TICK_COUNT + 1} är samma ettfel, och det är värt att
reda ut här, där du kan räkna flankerna med ögat, snarare än senare vid 50~MHz.

\exercise{Fyra frågor om modulen}{Resonemang}
\label{c7:ex:questions}
Svara på var och en av frågorna nedan i en eller två meningar, med hänvisning till modulen
\code{timer} i \secref{c7:sec:module}:

\xpart{a} Varför är \code{timeout} en encykelspuls snarare än en signal som ligger kvar hög när
målvärdet väl är nått?

\xpart{b} Medan en löpande timers \code{enable}-ingång hålls låg, vad händer med dess interna
räknare från en klockflank till nästa: nollställs den, pausar den, eller räknar den vidare?

\xpart{c} Varför tar \code{timer} \code{reset_s2_n} (den synkroniserade reset-signalen) snarare än
den råa \code{reset_n} som sin reset-ingång?

\xpart{d} CircuitVerse-timern i Övning~\ref{c7:ex:byhand} behövde kopplas om för att målvärdet
skulle ändras, medan VHDL-modulen inte behöver ändras alls. Vad ersätter omkopplingen, och när låses
dess värde fast?

% ------------------------------------------------------------------------------------------
\exerciseset{Timers i VHDL}

\exercise{Tre tickantal}{Resonemang}
\label{c7:ex:ticks}
DE0-CV-kortets systemklocka är \code{50 MHz}. Räkna för varje mål nedan ut den
\code{TICK_COUNT}-generic du skulle skicka in till modulen \code{timer}:

\xpart{a} En \code{timeout}-puls en gång per sekund.

\xpart{b} En \code{timeout}-puls var \code{250} ms.

\xpart{c} En \code{timeout}-puls tjugo gånger per sekund (\code{20 Hz}).

\note[Tips]Sambandet mellan en period och dess tickantal är
$\mathit{TICK\_COUNT} = \mathit{sekunder} \times 50\,000\,000 - 1$. Termen $-1$ är samma ettfel som
Övning~\ref{c7:ex:byhand} kretsade kring: räknaren löper \code{0} till och med \code{TICK_COUNT}, så
en period är \code{TICK_COUNT + 1} klockcykler, inte \code{TICK_COUNT}.

\exercise{\code{timer}}{VHDL}
\label{c7:ex:timer}
Skriv själva modulen \code{timer}, som en entitet vid namn \code{timer}.

Den byggs live under passet och är tryckt i sin helhet i \secref{c7:sec:module}, så det här är ingen
gåta. Skriv den ändå, utan att kopiera: varje design härifrån till bokens slut instansierar den här
modulen, båda capstones inräknade, och det är värt att ha byggt det du strax ska återanvända sex
gånger.

Entiteten har en generic:

\begin{booktable}{@{}p{6em}p{4.4em}p{6em}L@{}}
  \thead{Generic} & \thead{Typ} & \thead{Förval} & \thead{Betydelse} \\
  \midrule
  \code{TICK_COUNT} & \code{natural} & \code{50_000_000} & Antal tick att räkna innan en timeout.
    En sekund vid DE0-CV:ns \code{50 MHz}-klocka, avrundat: enligt Övning~\ref{c7:ex:ticks}:s regel
    är det exakta värdet \code{49_999_999}, och 20~ns fel på en sekund är inte värt den mindre
    läsbara konstanten. \\
\end{booktable}

och dessa portar:

\begin{booktable}{@{}p{5.6em}p{3.4em}p{5.2em}L@{}}
  \thead{Port} & \thead{Riktn.} & \thead{Typ} & \thead{Betydelse} \\
  \midrule
  \code{clock} & in & \code{std_logic} & Systemklocka. \\
  \code{reset_s2_n} & in & \code{std_logic} & Aktiv låg, redan synkroniserad reset. Asynkron: den
    nollställer räknaren och \code{timeout} utan att någon klockflank är inblandad. \\
  \code{enable} & in & \code{std_logic} & Timerns enable. Medan den är låg \textbf{håller} räknaren
    sitt värde, och fortsätter därifrån i stället för att börja om. \\
  \code{timeout} & out & \code{std_logic} & Encykelspuls, som slår till var \code{TICK_COUNT + 1}
    cykler medan timern är aktiverad. \\
\end{booktable}

Arkitekturen ska:
\begin{itemize}
  \item Deklarera den interna räknaren som \code{natural range 0 to TICK_COUNT}, så att registret
    dimensioneras av genericen snarare än lämnas obegränsat.
  \item Använda en enda process, känslig för \code{clock} och \code{reset_s2_n}, enligt mallen för
    klockad process från \secref{c3:sec:template}.
  \item Nollställa både räknaren och \code{timeout} vid reset, utanför den klockade grenen.
  \item Driva \code{timeout} låg villkorslöst högst upp i den klockade grenen, och hög bara vid den
    flank som nollställer räknaren. Det mönstret, ett förval följt av en överskrivning, är vad som
    gör pulsen exakt en cykel bred snarare än två.
\end{itemize}

Svara sedan, i en mening var:
\begin{itemize}
  \item \code{enable = '0'} måste \textbf{pausa} räknaren, inte nollställa den. Anta att du
    nollställde den i stället. Förutsäg vad referenstestbänken rapporterar, och vilken av de två
    egenskaperna ovan den mäter när den gör det. Kör den sedan och se. (\Chapref{c8:ch}:s
    \code{fsm_led} hänger på pausen, vilket är varför testbänken bryr sig.)
  \item Räknaren räknar \code{0} till och med \code{TICK_COUNT}, så en period är
    \code{TICK_COUNT + 1} cykler. Var i din kod är det där ”+ 1” faktiskt skrivet? Det står inte som
    en literal någonstans.
  \item \Secref{c7:sec:byhand}:s handritade \code{timeout} är en kombinatorisk avkodning av
    räkningen; din är registrerad. Båda har samma period. Vilken cykel slår var och en till på, och
    varför spelar skillnaden roll för något nedströms som samplar \code{timeout}?
\end{itemize}

\modulefig[0.62]{lectures/L07/appendix/images/timer.png}

\begin{selfcheck}
\textbf{Självkontroll.} Döp din entitet till \code{timer}, med genericen \code{TICK_COUNT}
(\code{natural}), ingångarna \code{clock}, \code{reset_s2_n}, \code{enable} och utgången
\code{timeout}, deklarerade i den ordningen; dess testbänk ligger i \file{exercises/timer/}. Den kör
tre hela perioder, så en timer som bara slår till en gång passerar inte, och den pausar mitt i
räkningen för att bekräfta att räkningen fortsatte i stället för att börja om.

\textbf{Spara den här filen.} Övning~\ref{c7:ex:blinker} nedan återanvänder den, det gör
\code{walking_led} i Övning~\ref{c7:ex:walking} också, och likaså båda \chapref{c8:ch}:s capstones.
\end{selfcheck}

\exercise{\code{blinker}}{VHDL}
\label{c7:ex:blinker}
Skriv en entitet vid namn \code{blinker} som blinkar en lysdiod på och av en gång per sekund, genom
att återanvända modulen \code{timer} du just skrev snarare än att räkna klockcykler själv.

Entiteten har en generic:

\begin{booktable}{@{}p{6em}p{4.4em}p{6em}L@{}}
  \thead{Generic} & \thead{Typ} & \thead{Förval} & \thead{Betydelse} \\
  \midrule
  \code{TICK_COUNT} & \code{natural} & \code{25_000_000} & En halv sekund vid DE0-CV:ns
    \code{50 MHz}-klocka. Exponera det som en generic i stället för att baka in antalet, så att en
    testbänk kan korta blinket; FPGA-bygget använder bara standardvärdet. \\
\end{booktable}

och dessa portar:

\begin{booktable}{@{}p{5em}p{3.4em}p{5.2em}L@{}}
  \thead{Port} & \thead{Riktn.} & \thead{Typ} & \thead{Betydelse} \\
  \midrule
  \code{clock} & in & \code{std_logic} & Systemklocka. \\
  \code{reset_n} & in & \code{std_logic} & Aktiv låg reset, rå från omvärlden. Asynkron: aktivera
    den och \code{led} slocknar omedelbart, utan att vänta på någon klockflank. \\
  \code{led} & out & \code{std_logic} & Nollställd till \code{'0'} medan reset-signalen är aktiv,
    och lämnad där till den första timeouten efter att reset-signalen släppts. \\
\end{booktable}

Lägg märke till att den här entiteten tar den \textbf{råa} \code{reset_n} snarare än en redan
synkroniserad \code{reset_s2_n}. Att synkronisera den är en del av övningen, och det är vad
\secref{c4:sec:resetsync} menar med ”exakt en modul i en design är undantaget”: här är den modulen
\code{blinker} själv.

Arkitekturen ska:
\begin{itemize}
  \item Instansiera din \code{reset_sync} från Övning~\ref{c4:ex:resetsync}, och använda dess
    utgång \code{reset_s2_n} som reset för allt inuti.
  \item Instansiera modulen \code{timer} från \secref{c7:sec:module}, skicka din egen generic
    \code{TICK_COUNT} rakt igenom till timerns \code{TICK_COUNT}, och hålla dess
    \code{enable}-ingång hög.
  \item Växla \code{led} i en synkron process varje gång timerns \code{timeout}-puls slår till.
\end{itemize}

Implementera inte räknelogiken, eller synkroniseraren, på nytt själv. Återanvänd \code{timer} och
\code{reset_sync} oförändrade, via en \code{generic map} och en \code{port map}, enligt
kompositionsmönstret i \secref{c7:sec:compose}. Förklara varför en timerperiod på en \textbf{halv}
sekund ger en blinkcykel på \textbf{en} sekund (en hel period av på och sedan av).

Svara sedan, i en mening var:
\begin{itemize}
  \item Lysdioden måste vara släckt under reset, och testbänken kontrollerar att den slocknar utan
    någon klockflank däremellan. Vilken halva av mönstret aktivera-asynkront\--släpp-synkront
    kontrollerar det, och vilken halva kontrollerar det inte?
  \item Din \code{timer} resettas av \code{reset_s2_n}, två flanker efter att \code{reset_n}
    släppts. Vad skulle gå fel om du matade den med den råa \code{reset_n} i stället? Namnge felet,
    inte bara regeln.
\end{itemize}

\modulefig[0.56]{lectures/L07/appendix/images/blinker.png}

\begin{selfcheck}
\textbf{Självkontroll.} Döp din entitet till \code{blinker}, med genericen \code{TICK_COUNT}
(\code{natural}), ingångarna \code{clock}, \code{reset_n} och utgången \code{led}, deklarerade i den
ordningen; dess testbänk ligger i \file{exercises/blinker/}. \textbf{Kopiera in dina egna
\file{reset_sync.vhd} och \file{timer.vhd} i den katalogen} innan du bygger; ingen av dem delas ut
där, eftersom du skrivit båda:

\begin{shell}
cd lectures/L07/exercises/blinker
cp ../../../L04/exercises/reset_sync/reset_sync.vhd .   # the two you wrote yourself
cp ../timer/timer.vhd .
ghdl -a --std=93 reset_sync.vhd timer.vhd blinker.vhd blinker_tb.vhd
\end{shell}

Testbänken skriver över \code{TICK_COUNT} med ett litet värde och kontrollerar att lysdioden växlar
om och om igen, ligger kvar exakt en timerperiod mellan växlingarna, och att den första växlingen
kommer sent med din synkroniserares två flanker.
\end{selfcheck}

% ------------------------------------------------------------------------------------------
\exerciseset{Den vandrande lysdioden}

En enda tänd lysdiod som vandrar längs en rad, en position per timertick, startad och stoppad med en
knapp. Den byggs live under passet och är beskriven i \secref{c7:sec:walking}.

Det här är utdelningen från de fyra senaste kapitlen, och den skriver nästan ingen egen logik: tre
moduler, två av dem oförändrade från \chapref{c4:ch} och en skriven här, komponerade.
Övning~\ref{c7:ex:walking} bygger den, och Övning~\ref{c7:ex:variants} och \ref{c7:ex:direction}
ändrar det du byggt.

\exercise{\code{walking_led}}{VHDL}
\label{c7:ex:walking}
Skriv en entitet vid namn \code{walking_led}.

Entiteten har två generics:

\begin{booktable}{@{}p{6em}p{4.4em}p{6em}L@{}}
  \thead{Generic} & \thead{Typ} & \thead{Förval} & \thead{Betydelse} \\
  \midrule
  \code{LED_COUNT} & \code{natural} & \code{8} & Hur många lysdioder biten vandrar längs. Raden är
    så här bred, och vandringen slår över efter så här många steg. \\
  \code{TICK_COUNT} & \code{natural} & \code{25_000_000} & Skickas rakt igenom till din
    \code{timer}, så att biten flyttar en position per timerperiod: ett steg per halv sekund vid
    DE0-CV:ns \code{50 MHz}-klocka. \\
\end{booktable}

och dessa portar:

\begin{booktable}{@{}p{5em}p{3.4em}p{11.4em}L@{}}
  \thead{Port} & \thead{Riktn.} & \thead{Typ} & \thead{Betydelse} \\
  \midrule
  \code{clock} & in & \code{std_logic} & Systemklocka. \\
  \code{reset_n} & in & \code{std_logic} & Aktiv låg reset, rå från omvärlden. Att synkronisera den
    är ditt jobb, precis som den var i \code{blinker}. \\
  \code{button_n} & in & \code{std_logic} & Aktiv låg tryckknapp, rå. Varje tryck startar eller
    stoppar vandringen. \\
  \code{led} & out & \code{std_logic_vector(LED_COUNT-1 downto 0)} & Lysdiodsraden. Exakt en bit är
    tänd åt gången; vid reset är det bit \code{0}. \\
\end{booktable}

Arkitekturen komponerar tre moduler du redan skrivit, och lägger till två processer:
\begin{itemize}
  \item Instansiera din \textbf{\code{reset_sync}}, och använd dess \code{reset_s2_n} som reset för
    allt annat i designen, de två instanserna nedan inräknade.
  \item Instansiera din \textbf{\code{button_sync}} med \code{generic map(1)}, för en encykelspuls
    per tryck. Dess knappportar är vektorer med \code{COUNT} bitar medan \code{button_n} här är en
    skalär, så överbrygga bredden med en \code{std_logic_vector(0 downto 0)} med ett element i
    vardera riktningen. \Secref{c7:sec:walking} visar det.
  \item Instansiera din \textbf{\code{timer}} med din egen \code{TICK_COUNT}, driv dess
    \code{enable} från skiftflaggan nedan och ta dess \code{timeout} som skiftticket.
  \item En process som växlar \textbf{\code{shift_enable}} vid varje puls från knappen, nollställd
    till \code{'0'} vid reset. Vandringen börjar därför stoppad.
  \item En process som håller det \textbf{skiftregister} som driver \code{led}: vid reset, ladda en
    enda tänd bit på position \code{0} och nollställ resten; vid varje \code{timeout} från timern,
    rotera ett steg mot MSB och låt den översta biten slå över tillbaka in i bit \code{0}. Det är
    ett PISO-register (\secref{c6:sec:piso}) med sin egen seriella utgång kopplad tillbaka till sin
    seriella ingång, vilket är det som gör en engångsutskiftning till en oändlig vandring.
\end{itemize}

Svara sedan, i en mening var:
\begin{itemize}
  \item \code{shift_enable} driver timerns \code{enable}, så att stoppa vandringen stoppar timern
    också. Övning~\ref{c7:ex:timer} fick dig att visa att \code{enable} \textbf{pausar} räkningen
    snarare än nollställer den. Vad ser en betraktare på kortet, vid det tryck som startar om
    vandringen, om din timer nollställer i stället?
  \item \code{button_sync} instansieras med \code{COUNT = 1} här och med \code{COUNT = 2} i
    \chapref{c8:ch}:s \code{fsm_led}, ur samma fil utan någon ändring. Vad hade du behövt skriva i
    stället om bredden var fastlagd i modulen snarare än inskickad som en generic?
\end{itemize}

\modulefig[0.7]{lectures/L07/appendix/images/walking_led.png}

\begin{selfcheck}
\textbf{Självkontroll.} Döp din entitet till \code{walking_led}, med generics \code{LED_COUNT} och
\code{TICK_COUNT} (båda \code{natural}), ingångarna \code{clock}, \code{reset_n}, \code{button_n}
och utgången \code{led}, deklarerade i den ordningen; dess testbänk ligger i
\file{exercises/walking_led/}. Den skriver över båda generics (\code{LED_COUNT = 4},
\code{TICK_COUNT = 3}) så att ett helt varv tar en handfull klockcykler snarare än två sekunder, och
den kontrollerar resetmönstret, att ingenting rör sig före det första trycket, och att antalet
skiftningar över ett fast fönster stämmer med den takt timern sätter. Kopiera först in de tre
moduler den komponerar; kommandot står i \secref{c7:sec:walking}.
\end{selfcheck}

\exercise{Att ändra takt och bredd}{Simulering}
\label{c7:ex:variants}
Gör följande ändringar en i taget, och kör om testbänken efter var och en. Del \textbf{a)} och
\textbf{b)} ändrar bara testbänkens konstanter \textemdash{} att de inte kräver någon ändring alls i
din \file{walking_led.vhd} är hela poängen med dem. Del \textbf{c)} är den som rör modulen.

\xpart{a} Ändra vandringshastigheten:
\begin{itemize}
  \item Ändra konstanten \code{TICK_COUNT} i \textbf{testbänken} från \code{3} till \code{6}.
  \item Den ska fortfarande passera, den här gången för att testbänken anpassar sig med dig: dess
    observationsfönster är fasta \code{PULSES * 6} cykler, och den härleder det antal skiftningar
    den förväntar sig ur \code{TICK_COUNT}, en skiftning per \code{TICK_COUNT + 1} cykler.
  \item Hitta nu var den anpassningen slutar betyda något. Fortsätt höja \code{TICK_COUNT} förbi
    fönstrets längd och räkna ut vad det förväntade antalet blir, och vad kontrollen då faktiskt
    bevisar. Ett test som passerar av fel skäl är värt att kunna känna igen.
  \item Förklara sedan varför ett ändrat tickantal inte kan göra sönder \emph{designen}, bara sakta
    ner den.
\end{itemize}

\xpart{b} Ändra bredden: ändra testbänkens konstant \code{LED_COUNT} från \code{4} till \code{8}.
Den ska fortfarande passera, utan någon ändring alls i din \file{walking_led.vhd}. Förklara hur
modulen hade behövt se ut för att det här skulle ha krävt en ändring.

\xpart{c} Ladda två tända bitar på var sin sida av registret vid reset (t.ex. \code{"10001000"} för
\code{LED_COUNT = 8}), och låt dem vandra i takt:
\begin{itemize}
  \item Peka ut den enda rad i din \file{walking_led.vhd} du behöver ändra.
  \item \textbf{Förutsäg} vad referenstestbänken kommer att rapportera, innan du kör någonting.
  \item Kör den sedan och se. Förklara resultatet: är designen fel, eller kontrollerar testbänken
    något som inte längre är specifikationen?
\end{itemize}

\note[Tips]Både \code{LED_COUNT} och \code{TICK_COUNT} är generics, så del \textbf{a)} och
\textbf{b)} kräver ingen ändring alls i modulens kropp. Det är hela poängen med en generic: en
modul, många storlekar, inga ändringar.

\exercise{En andra knapp som vänder riktningen}{VHDL}
\label{c7:ex:direction}
Den nuvarande designen vandrar bara åt ett håll (mot MSB, där den översta biten slår över tillbaka
till botten). Ändra din \file{walking_led.vhd} så att en \textbf{andra} knapp vänder
vandringsriktningen.

\xpart{a} Lägg till en andra, synkroniserad knappingång: dra den genom samma
\code{button_sync}-instans, och skicka \code{2} till dess generic \code{COUNT}
(\code{generic map(2)}, positionsvis, som överallt annars i den här boken).

\xpart{b} Lägg till en signal \code{direction}, och växla den vid den nya knappens flank, precis så
som \code{shift_enable} växlas av den första knappen.

\xpart{c} Gör rotationen i \code{SHIFT_PROCESS} beroende av \code{direction}. När
\code{direction = '0'}, rotera mot MSB (bokens konvention från \secref{c6:sec:shift}, och det du
skrev i Övning~\ref{c7:ex:walking}):

\begin{vhdlcode}
shift_reg <= shift_reg(LED_COUNT-2 downto 0) & shift_reg(LED_COUNT-1);
\end{vhdlcode}

När \code{direction = '1'}, rotera åt andra hållet, mot bit 0:

\begin{vhdlcode}
shift_reg <= shift_reg(0) & shift_reg(LED_COUNT-1 downto 1);
\end{vhdlcode}

Det här är det enda stället i boken där ett skiftregister medvetet går emot konventionen: att vända
riktningen är hela poängen med övningen. Lägg märke till att det andra uttrycket är precis den
spegelbild \secref{c6:sec:shift} varnar dig för att hålla utkik efter i andras kod.

\xpart{d} Din entitets portar har nu ändrats: \code{button_n} är en tvåbitarsvektor snarare än en
enda bit.
\begin{itemize}
  \item Förklara varför \file{walking_led_tb.vhd} inte längre går att köra mot din version, och
    exakt vad GHDL kommer att klaga på när du försöker.
  \item Det här är den vanliga kostnaden för att ändra ett gränssnitt, och den är värd att känna en
    gång: en testbänk är skriven mot en bestämd entitet, och att bredda en port gör sönder varje
    instansiering av den, också de du inte skrivit själv.
  \item Beskriv, i löpande text, den följd av tryck du skulle använda för att övertyga dig själv om
    att designen fungerar, och vad du skulle förvänta dig att lysdioderna gör vid varje steg. Den
    följden är precis vad en testbänk för den nya entiteten skulle behöva driva.
\end{itemize}
